\documentclass[10pt,aspectratio=169]{beamer}
\usetheme{CambridgeUS}
\usecolortheme{dolphin}
\usepackage{amsmath,amssymb,amsthm}

\title{Theorem, Proof, And Consequences}
\author{Speaker Name}
\date{\today}

\begin{document}
\begin{frame}\titlepage\end{frame}
\begin{frame}{Roadmap}\tableofcontents\end{frame}

\section{Setup}
\begin{frame}{Objects And Assumptions}
  \begin{definition}
    Define the main objects and their scope.
  \end{definition}
  \begin{block}{Standing assumptions}
    List assumptions that must not disappear from later slides.
  \end{block}
\end{frame}

\section{Theorem}
\begin{frame}{Main Theorem}
  \begin{theorem}[Main result]
    State the theorem exactly as supported by the source.
  \end{theorem}
  \begin{alertblock}{Meaning risk}
    Record any condition that is simplified in speech but not omitted.
  \end{alertblock}
\end{frame}

\section{Proof}
\begin{frame}{Proof Roadmap}
  \begin{enumerate}
    \item Normalization or reduction.
    \item Key lemma.
    \item Final assembly.
  \end{enumerate}
\end{frame}

\begin{frame}{Key Lemma}
  \begin{lemma}
    State the lemma with the exact hypotheses needed by the theorem.
  \end{lemma}
  \begin{proof}[Proof idea]
    Give only the idea; put full details in backup unless the talk requires them.
  \end{proof}
\end{frame}

\section{Consequences}
\begin{frame}{Consequences And Sharpness}
  \begin{itemize}
    \item Corollary or example.
    \item Limitation or non-example.
    \item Open problem.
  \end{itemize}
\end{frame}
\end{document}
